\section{Results}

The proposed MIND-BLOOM framework was evaluated on the Bangladesh PPD dataset ($n=769$ after preprocessing) using an 80-20 stratified train-test split. All models were trained on SMOTENC-balanced training data (615 samples balanced to 271 per class) and evaluated on the original, unbalanced test set (154 samples) to assess real-world performance.

\subsection{Model Performance Overview}

Table \ref{tab:model_performance} presents comprehensive performance metrics for all evaluated models. Random Forest emerged as the best-performing model across key metrics, followed closely by XGBoost and CatBoost.

\begin{table*}[htbp]
\centering
\caption{Comprehensive Model Performance Comparison}
\label{tab:model_performance}
\begin{tabular}{lcccccc}
\toprule
\textbf{Model} & \textbf{Accuracy} & \textbf{Sensitivity} & \textbf{Specificity} & \textbf{F1-Score} & \textbf{AUC-ROC} & \textbf{MCC} \\
\midrule
\textbf{Random Forest} & \textbf{0.8333} & \textbf{0.8333} & \textbf{0.9167} & \textbf{0.8328} & \textbf{0.9469} & \textbf{0.7502} \\
XGBoost (Optimized) & 0.8284 & 0.8284 & 0.9142 & 0.8274 & 0.9406 & 0.7434 \\
CatBoost & 0.8186 & 0.8186 & 0.9093 & 0.8185 & 0.9554 & 0.7371 \\
LightGBM & 0.8186 & 0.8186 & 0.9093 & 0.8179 & 0.9406 & 0.7361 \\
Logistic Regression & 0.8137 & 0.8137 & 0.9069 & 0.8129 & 0.9309 & 0.7209 \\
TabNet & 0.7843 & 0.7843 & 0.8922 & 0.7861 & 0.9082 & 0.6792 \\
K-Nearest Neighbors & 0.7451 & 0.7451 & 0.8725 & 0.7478 & 0.8962 & 0.6238 \\
Decision Tree & 0.6961 & 0.6961 & 0.8480 & 0.6966 & 0.8124 & 0.5442 \\
SVM (RBF) & 0.3873 & 0.3873 & 0.6936 & 0.3293 & 0.4886 & 0.1079 \\
\bottomrule
\end{tabular}
\end{table*}

\subsection{Clinical Performance Interpretation}

\textbf{Random Forest achieved a sensitivity of 83.3\%}, correctly identifying 83 out of 100 mothers with actual postpartum depression, with an \textbf{AUC-ROC of 0.95} and \textbf{F1-score of 0.83}. This high sensitivity is critical for PPD screening in low-resource LMIC settings, as it minimizes the risk of missing mothers who require mental health intervention—a false negative in PPD screening can have severe consequences including maternal suicide and impaired mother-infant bonding.

The model also achieved a specificity of 91.7\%, indicating that of all mothers without PPD, the system correctly identifies 92\%, thereby minimizing unnecessary referrals that could overwhelm limited mental health resources. The AUC-ROC of 0.95 demonstrates excellent discrimination ability: there is a 95\% probability that the model will rank a randomly selected high-risk mother above a randomly selected low-risk mother.

\subsection{Advanced Evaluation Metrics}

Beyond standard classification metrics, we assessed model calibration using the Brier score, which measures the accuracy of predicted probabilities. Random Forest achieved a Brier score of 0.102 (where lower is better and scores below 0.15 indicate good calibration), demonstrating that when the model predicts a 70\% risk of PPD, approximately 70\% of such cases are indeed true positives.

The Matthews Correlation Coefficient (MCC) of 0.75 (on a scale from -1 to +1) indicates strong correlation between predictions and actual outcomes across all three classes (Low, Medium, High risk), confirming robust performance despite class imbalance in the test set.

Decision curve analysis (Fig. \ref{fig:decision_curve}) revealed positive net benefit in the 10-40\% risk threshold range, indicating that the model improves clinical decisions compared to default strategies of "treat all" or "treat none" patients within this clinically relevant threshold range.

\subsection{Model Calibration and Fairness}

Calibration curves (Fig. \ref{fig:calibration}) demonstrate excellent agreement between predicted probabilities and observed frequencies across all risk levels. The model maintains consistent performance across demographic subgroups, with accuracy variations of less than 3\% across education levels, residence types, family structures, and age groups, ensuring equitable performance across diverse populations.

\subsection{Comparison with State-of-the-Art}

Table \ref{tab:comparison_literature} contextualizes our results within the existing PPD prediction literature. Our Random Forest model achieved an AUC-ROC of 0.95, approaching the performance of Chinese tertiary hospital models (Huang et al., 2025: AUC 0.955) \cite{Huang2025}, while relying solely on survey-based features without laboratory tests or clinical biomarkers. This demonstrates that high-accuracy PPD prediction is achievable in resource-constrained LMIC settings using accessible screening tools.

\begin{table}[htbp]
\centering
\caption{Performance Comparison with Existing PPD Prediction Studies}
\label{tab:comparison_literature}
\begin{tabular}{p{0.12\linewidth}p{0.18\linewidth}p{0.15\linewidth}p{0.12\linewidth}p{0.22\linewidth}}
\toprule
\textbf{Study} & \textbf{Setting} & \textbf{Method} & \textbf{AUC-ROC} & \textbf{Features} \\
\midrule
\multicolumn{5}{c}{\textit{International Studies (High-Resource Settings)}} \\
\midrule
Huang et al. (2025) & China Hospital & XGBoost + Boruta & 0.955 & Psychosocial + Labs \\
Zhang et al. (2025) & China Hospital & XGBoost + SHAP & 0.85 & EMR + Labs \\
Qi et al. (2025) & China Hospital & Logistic Reg. & 0.80-0.86 & Psychosocial \\
\midrule
\multicolumn{5}{c}{\textit{Bangladesh Studies (LMIC Settings)}} \\
\midrule
Zohora (2025) & Bangladesh & Random Forest & 0.86-0.98 & Clinical (private data) \\
Islam et al. (2025) & Bangladesh & Logistic Reg. & N/R & Socio-demographic \\
\midrule
\textbf{Current Study} & \textbf{Bangladesh} & \textbf{RF + CatBoost} & \textbf{0.95} & \textbf{Survey only (no labs)} \\
\bottomrule
\end{tabular}
\end{table}

Compared to existing Bangladesh-specific studies, our work represents the first comprehensive machine learning approach on an open dataset with rigorous evaluation protocols. While Zohora (2025) reported AUC values of 0.86-0.98 using Random Forest \cite{Zohora2025}, their dataset remains unpublished and focused on suicidal ideation rather than PPD classification. Islam et al. (2025) \cite{Islam2025} and Al-Mamun et al. (2024) \cite{AlMamun2024} identified important risk factors through traditional regression but did not develop predictive models.

\subsection{Feature Importance and Interpretability}

Feature importance analysis via SHAP values and TabNet attention mechanisms (Fig. \ref{fig:feature_importance}) revealed that PHQ-9 score, depression history, social support index, and cumulative risk score were the most influential predictors. TabNet attention analysis showed PHQ-9 Score (importance: 0.366), PHQ-9 Result (0.131), and Age of immediate older children (0.073) as the strongest predictors. The newly engineered \textit{low\_support\_flag} feature demonstrated meaningful contribution (importance: 0.018), validating our domain-informed feature engineering approach. This aligns with established PPD literature emphasizing the role of prenatal depression, inadequate social support, and cumulative psychosocial stressors \cite{Islam2025,AlMamun2024}.

\begin{figure}[htbp]
    \centering
    \includegraphics[width=0.9\linewidth]{tabnet_feature_importance.png}
    \caption{TabNet global feature importance via attention weights. PHQ-9 Score, PHQ-9 Result, and Age of older children emerge as the strongest predictors, consistent with PPD risk factor literature.}
    \label{fig:feature_importance}
\end{figure}

TabNet's attention masks provided patient-level explanations, showing which features the model focused on for each prediction. The interpretable feature selection mechanism, combined with SHAP analysis on traditional models, enables clinicians to understand and trust model predictions, facilitating clinical adoption in LMIC settings.

\subsection{Gradient Boosting Model Performance}

State-of-the-art gradient boosting frameworks CatBoost and LightGBM achieved accuracies of 81.86\%, with AUC-ROC values of 0.96 and 0.94, respectively. CatBoost's native handling of categorical features via ordered target encoding proved particularly effective for this survey-based dataset, achieving the highest AUC-ROC (0.9554) among all models. These models represent viable alternatives for deployment scenarios where PyTorch dependencies are unavailable, as they offer faster inference times and require less computational infrastructure.

\subsection{Deep Learning Model Performance}

TabNet, the attention-based deep learning architecture specifically designed for tabular data, achieved 78.43\% accuracy with an AUC-ROC of 0.91 and F1-score of 0.79. While TabNet's performance was lower than traditional ensemble methods for this dataset size (769 samples), it provided valuable built-in interpretability through attention mechanisms. The attention weights revealed clinically meaningful feature selection patterns, suggesting TabNet's utility in scenarios where model explainability is prioritized over marginal accuracy gains.

\subsection{Model Selection for Deployment}

Based on comprehensive evaluation, \textbf{Random Forest is recommended for primary deployment} due to:
\begin{itemize}
    \item Highest overall performance (accuracy: 83.33\%, AUC: 0.95, F1: 0.83)
    \item Best sensitivity (83.3\%) - critical for screening applications
    \item Excellent specificity (91.7\%) - minimizes false positives
    \item Highest MCC (0.75) - robust multiclass performance
    \item Well-established reliability in medical applications
    \item Fast inference without deep learning dependencies
\end{itemize}

\textbf{CatBoost is recommended as a strong alternative} for scenarios requiring:
\begin{itemize}
    \item Native categorical feature handling (no preprocessing needed)
    \item Highest AUC-ROC (0.9554) for risk stratification
    \item Comparable accuracy (81.86\%) with faster training
    \item Excellent feature importance interpretability
\end{itemize}

\textbf{TabNet serves as an interpretability-focused option}, offering attention-based explanations at the cost of ~5\% accuracy compared to Random Forest, suitable for research and clinical validation scenarios where understanding feature interactions is paramount.

\begin{figure}[htbp]
    \centering
    \includegraphics[width=0.85\linewidth]{decision_curve_analysis.png}
    \caption{Decision curve analysis showing positive net benefit in the 10-40\% risk threshold range.}
    \label{fig:decision_curve}
\end{figure}

\begin{figure}[htbp]
    \centering
    \includegraphics[width=0.85\linewidth]{calibration_curves.png}
    \caption{Calibration curves demonstrating excellent agreement between predicted probabilities and observed frequencies.}
    \label{fig:calibration}
\end{figure}

